\chapter{Groups, Quotients, and Homomorphisms}\label{chap:groups-foundations}

\section{Groups, Subgroups, and Generated Subgroups}

Groups record a collection of reversible symmetries or operations.  The
axioms are deliberately sparse: associativity makes an unparenthesized product
meaningful, the identity records ``doing nothing,'' and inverses record the
possibility of undoing an operation.  The familiar examples \(\mathbb Z\),
\(\mathbb Z/n\mathbb Z\), the symmetric groups \(S_n\), and the dihedral groups
\(D_n\) will recur throughout Part~I.

\begin{definition}
A \emph{group} is a set \(G\) equipped with a binary operation
\(G\times G\to G\), \((x,y)\mapsto xy\), satisfying the following axioms.
\begin{enumerate}
\setlength{\itemsep}{7pt}
\item \textbf{Associativity:} for all \(x,y,z\in G\),
      \[(xy)z=x(yz).\]
\item \textbf{Identity:} there is an element \(1\in G\) such that, for every
      \(x\in G\),
      \[1x=x=x1.\]
\item \textbf{Inverses:} for every \(x\in G\), there is an element
      \(x^{-1}\in G\) such that
      \[xx^{-1}=1=x^{-1}x.\]
\end{enumerate}
A nonempty subset \(H\subseteq G\) is a \emph{subgroup} when the same
operation makes \(H\) into a group; equivalently, it is closed under products
and inverses.
\end{definition}

Associativity permits \(g_1g_2\cdots g_n\) without specifying parentheses.

\begin{example}[Additive Number Systems]
The integers \((\mathbb Z,+)\) form a group: associativity is inherited from
ordinary addition, the identity is \(0\), and the inverse of \(m\) is \(-m\).
The same verification gives additive groups \((\mathbb Q,+)\) and
\((\mathbb R,+)\).  These examples explain the word ``inverse'' in an additive
setting: it is an additive inverse, not a reciprocal.  The natural numbers
\((\mathbb N,+)\) are not a group, because a positive integer has no additive
inverse in \(\mathbb N\).
\end{example}

\begin{example}[Multiplicative Number Systems]
The nonzero rational numbers \(\mathbb Q^\times=\mathbb Q\setminus\{0\}\)
form a group under multiplication.  The identity is \(1\), and the inverse of
\(a/b\) is \(b/a\).  Likewise \(\mathbb R^\times\) is a group.  The zero
rational or real number must be removed: it has no multiplicative inverse.
This small change illustrates a common algebraic move---restrict to the
elements for which an operation is reversible.
\end{example}

\begin{example}[Matrices]
For the familiar number systems \(F=\mathbb Q\) or \(\mathbb R\), the
invertible \(n\times n\) matrices form the group \(\operatorname{GL}_n(F)\).
Matrix multiplication is associative, the identity is \(I_n\), and membership
means precisely that a matrix inverse exists.  (The general notion of a field
will be developed later.)  For \(n\geq2\) this group is generally nonabelian:
two reversible linear transformations need not commute.  This is the first
indication that groups capture more than arithmetic.
\end{example}

\begin{proposition}[Elementary Group Identities]
The identity and inverses are unique.  Cancellation holds, and
\((xy)^{-1}=y^{-1}x^{-1}\).
\end{proposition}
\begin{proof}
If \(e,e'\) are identities, use first the right-identity property of \(e'\)
and then the left-identity property of \(e\): \(e=ee'=e'\).  If \(y,z\) both
invert \(x\), then associativity lets us insert parentheses as follows:
\[y=y1=y(xz)=(yx)z=1z=z.\]
For cancellation, if \(ax=ay\), multiply the equality on the left by
\(a^{-1}\):
\[x=1x=(a^{-1}a)x=a^{-1}(ax)=a^{-1}(ay)=(a^{-1}a)y=1y=y.\]
Right cancellation is similar.  Finally,
\[(xy)(y^{-1}x^{-1})=x(yy^{-1})x^{-1}=x1x^{-1}=1,\]
and the same calculation in the opposite order shows that \(y^{-1}x^{-1}\)
is the inverse of \(xy\).
\end{proof}

Thus a subgroup is not merely a subset: it is a collection of operations that
is itself closed under composition and reversal.  For example,
\(n\mathbb Z\leq\mathbb Z\), while the transpositions in \(S_3\) do not form a
subgroup because their product need not be a transposition.

We will use one kind of map before studying maps systematically.  A function
\(f:G\to H\) between groups is a \emph{homomorphism} if
\(f(xy)=f(x)f(y)\) for all \(x,y\in G\).  Its \emph{kernel} is the set of
elements sent to the identity of \(H\).  The structural consequences of these
definitions---images, normality, and quotient groups---are developed later in
this chapter.

\begin{proposition}[Subgroup Test]
A nonempty subset \(H\) of \(G\) is a subgroup if and only if
\(xy^{-1}\in H\) for all \(x,y\in H\).
\end{proposition}
\begin{proof}
If \(H\) is already a subgroup, then \(y^{-1}\in H\) and closure under
products gives \(xy^{-1}\in H\).  Conversely, assume the displayed condition.
Choose \(h\in H\), possible because \(H\) is nonempty.  Taking \(x=y=h\)
gives \(hh^{-1}=1\in H\).  Now take \(x=1\) and any \(y\in H\); the condition
gives \(y^{-1}\in H\).  Finally, if \(x,y\in H\), then \(y^{-1}\in H\), so
the condition applied to \(x\) and \(y^{-1}\) gives
\(x(y^{-1})^{-1}=xy\in H\).  Thus \(H\) contains the identity and is closed
under products and inverses.
\end{proof}

The subgroup test is useful because it replaces three separate closure checks
by a single expression.  It also anticipates a recurring algebraic pattern:
subobjects are often recognized by closure under a difference-like operation.

\begin{definition}
For a subset \(S\subseteq G\), the \emph{subgroup generated by \(S\)} is
\[
\langle S\rangle=\bigcap_{\substack{H\leq G\\S\subseteq H}}H.
\]
It is the smallest subgroup of \(G\) containing \(S\).  A group is
\emph{finitely generated} if \(G=\langle S\rangle\) for some finite \(S\).
\end{definition}

The intersection is a subgroup, and it is contained in every subgroup that
contains \(S\), which proves the asserted minimality.  Equivalently,
\(\langle S\rangle\) consists exactly of finite products of elements of \(S\)
and their inverses (including the empty product \(1\)).  This formulation is
often the practical one; the intersection formulation is the structural one.
The same ``smallest subobject containing a set'' construction will reappear for
ideals, modules, algebras, and field extensions.

\begin{example}
In the additive group \(\mathbb Z\), \(\langle6,10\rangle=2\mathbb Z\), since
\(2=2\cdot6-10\).  In \(S_3\), \((12)\) and \((23)\) generate \(S_3\).
\end{example}

\begin{proposition}[Finite Products from a Generating Set]
For \(S\subseteq G\), the subgroup \(\langle S\rangle\) consists exactly of
the finite products \(s_1^{\varepsilon_1}\cdots s_k^{\varepsilon_k}\), where
\(s_i\in S\) and \(\varepsilon_i\in\{1,-1\}\), together with the empty
product \(1\).
\end{proposition}
\begin{proof}
Let \(W\) be the set of these finite products in the already given group
\(G\).  Concatenation gives closure under products, and reversal with inverse
exponents gives closure under inverses.
Thus \(W\) is a subgroup containing \(S\), so minimality gives
\(\langle S\rangle\subseteq W\).  Every subgroup containing \(S\) contains
all these finite products, so its intersection contains \(W\).
\end{proof}

This two-containment argument will recur whenever we describe an object
generated by a specified set of elements.

\section{Cyclic Groups}

For \(g\in G\), \(\langle g\rangle=\{g^n:n\in\mathbb Z\}\) is the special
case \(S=\{g\}\), the \emph{cyclic subgroup generated by \(g\)}.  The order
\(|g|\) is the least positive \(n\) with \(g^n=1\), if it exists.

The cyclic group \(\langle g\rangle\) is the portion of \(G\) generated by one
operation.  Additively, \(\langle 1\rangle=\mathbb Z\); multiplicatively, the
rotations of a regular \(n\)-gon form \(C_n\).  The order of \(g\) measures the
first time repeated application returns to the identity.

\begin{proposition}
Every subgroup of a cyclic group is cyclic.  If \(G=\langle g\rangle\) has
order \(n\), then for each divisor \(d\mid n\), \(G\) has a unique subgroup of
order \(d\), namely \(\langle g^{n/d}\rangle\).
\end{proposition}
\begin{proof}
For a nontrivial subgroup \(H\), let \(r\) be the least positive integer with
\(g^r\in H\).  Given \(g^k\in H\), divide \(k\) by \(r\): write
\(k=qr+s\), \(0\leq s<r\).  Then \(g^s=g^k(g^r)^{-q}\in H\).  Minimality of
\(r\) forces \(s=0\), so every element of \(H\) is a power of \(g^r\), and
\(H=\langle g^r\rangle\).  In a finite cyclic group, \(g^a\) has order
\(n/\gcd(n,a)\).  Thus \(g^{n/d}\) has order \(d\).  If a subgroup has order
\(d\), write it as \(H=\langle g^r\rangle\).  Then
\(\gcd(n,r)=n/d\), so \(r=(n/d)u\) with \(\gcd(u,d)=1\).  Since
\(g^{n/d}\) has order \(d\), there are integers \(v,w\) with
\(uv+dw=1\).  Consequently
\[
 g^{n/d}=(g^r)^v(g^{n/d})^{dw}=(g^r)^v,
\]
because \((g^{n/d})^d=1\).  Hence
\(\langle g^{n/d}\rangle\leq\langle g^r\rangle\); the reverse inclusion is
immediate from \(g^r=(g^{n/d})^u\).  Thus
\(H=\langle g^{n/d}\rangle\), which proves uniqueness precisely.
\end{proof}

\begin{example}[The Lattice of \(C_{12}\)]
\leavevmode\\
If \(C_{12}=\langle g\rangle\), its subgroups are
\(\langle g\rangle,\langle g^2\rangle,\langle g^3\rangle,
\langle g^4\rangle,\langle g^6\rangle\), and \(\{1\}\), of orders
\(12,6,4,3,2,1\), respectively.  The divisors of \(12\) predict exactly this
subgroup lattice.  In contrast, the analogous subgroup problem for a general
finite group is substantially harder, which is why cyclic groups are a useful
first laboratory.
\end{example}

\begin{proposition}
The homomorphism \(\varphi_g:\mathbb Z\to G\), \(n\mapsto g^n\), has image
\(\langle g\rangle\), and
\[
 \ker\varphi_g=
 \begin{cases}
 \{0\},&\text{if \(g\) has infinite order},\\[2pt]
 n\mathbb Z,&\text{if \(|g|=n\)}.
 \end{cases}
\]
Thus an infinite cyclic group is isomorphic to \(\mathbb Z\), whereas a
cyclic group of order \(n\) is isomorphic to \(\mathbb Z/n\mathbb Z\).
\end{proposition}
\begin{proof}
For integers \(a,b\),
\(\varphi_g(a+b)=g^{a+b}=g^ag^b=\varphi_g(a)\varphi_g(b)\), so
\(\varphi_g\) is indeed a homomorphism.  Its image is precisely the set of
all integer powers of \(g\), namely \(\langle g\rangle\).
The kernel consists of precisely those integers \(a\)
for which \(g^a=1\).  If \(g\) has infinite order, this happens only for
\(a=0\).  If \(|g|=n\), division with remainder shows that \(g^a=1\) if and
only if \(n\mid a\), so the kernel is \(n\mathbb Z\).  In the infinite-order
case, \(\varphi_g\) is one-to-one: \(g^a=g^b\) implies
\(g^{a-b}=1\), hence \(a=b\).  It is onto \(\langle g\rangle\) by its
definition, and thus gives \(\mathbb Z\cong\langle g\rangle\).
If \(|g|=n\), the map
\[a+n\mathbb Z\longmapsto g^a\]
is well-defined because \(g^n=1\); it is a bijective homomorphism onto
\(\langle g\rangle\).  The First Isomorphism Theorem will later combine these
two direct arguments into the single statement
\(\langle g\rangle\cong\mathbb Z/\ker\varphi_g\).
\end{proof}

\begin{proposition}
If \(G=\langle g\rangle\) has order \(n\), then \(g^k\) is a generator if and
only if \(\gcd(k,n)=1\).
\end{proposition}
\begin{proof}
The order of \(g^k\) is \(n/\gcd(n,k)\).  Thus \(C_{12}\) has generators
\(g,g^5,g^7,g^{11}\).
\end{proof}

\section{Dihedral Groups}

\begin{definition}
The \emph{dihedral group} \(D_n\) is the full symmetry group of a regular
\(n\)-gon.  If \(r\) is rotation through one vertex and \(s\) a reflection,
then
\[D_n=\langle r,s\mid r^n=s^2=1,\ srs^{-1}=r^{-1}\rangle.\]
Many authors call this order-\(2n\) group \(D_{2n}\); these notes use \(D_n\).
\end{definition}

The relation \(sr=r^{-1}s\) lets us move every occurrence of \(s\) to the
left in a product.  Since \(s^2=1\) and \(r^n=1\), every product of rotations
and reflections therefore has the form \(r^i\) or \(sr^i\), with
\(0\leq i<n\).  These \(2n\) symmetries are distinct geometrically: the first
family preserves orientation and the second reverses it.  Hence
\(|D_n|=2n\) and \(\langle r\rangle\cong C_n\).  Notice that multiplying
\(srs^{-1}=r^{-1}\) on the right by \(s\) gives the useful rewriting rule
\(sr=r^{-1}s\): reflecting after rotating has the same effect as rotating in
the opposite direction after reflecting.

For the square, choose \(r\) to be the quarter-turn and \(s\) a reflection
in an axis through opposite vertices.  The complete list is
\[
 D_4=\{1,r,r^2,r^3,s,sr,sr^2,sr^3\}.
\]
The first four are rotations, and the last four are reflections.  In
particular, \(sr=r^{-1}s=r^3s\), whereas \(rs\) is a different reflection:
\(sr\ne rs\).  This is a concrete calculation of noncommutativity, not just
an abstract consequence of the presentation.

\begin{example}[Dihedral Groups as Permutation Groups]
Label the vertices of a regular \(n\)-gon by \(1,\ldots,n\).  Every symmetry
permutes these vertices, so it determines an element of \(S_n\).  The action
is injective: a geometric symmetry that fixes every vertex is the identity.
Thus \(D_n\) is concretely a subgroup of \(S_n\).  The later Cayley theorem
will generalize this phenomenon from dihedral groups to every finite group.
\end{example}

\section{Symmetric and Alternating Groups}

\begin{definition}
The \emph{symmetric group} \(S_n\) is the group of all bijections of the set
\(\{1,2,\ldots,n\}\), with composition as the operation.  A \(k\)-\emph{cycle}
\((a_1\ a_2\ \cdots\ a_k)\) sends \(a_i\) to \(a_{i+1}\) for
\(1\leq i<k\), sends \(a_k\) to \(a_1\), and fixes every other letter.
\end{definition}

Throughout this section, products of permutations are evaluated from right to
left: \(\sigma\tau\) means ``first apply \(\tau\), then apply \(\sigma\).''
Cycle notation records only the letters that move.  For example,
\((1\ 3\ 2)\in S_4\) sends \(1\mapsto3\), \(3\mapsto2\),
\(2\mapsto1\), and fixes \(4\).  A cycle may be cyclically rewritten without
changing the permutation: \((1\ 3\ 2)=(3\ 2\ 1)\).  Cycles with disjoint
supports commute because they act independently on every letter.  For a
worked noncommuting product,
\[
 (1\ 2\ 3)(1\ 2)=(1\ 3).
\]
Indeed, \(1\mapsto2\mapsto3\), \(3\mapsto3\mapsto1\), and \(2\) is fixed
by the composite.  Reading this calculation from the right is essential.

\begin{proposition}[Cycle Decomposition]
Every permutation in \(S_n\) is a product of disjoint cycles.  This
decomposition is unique up to reordering the cycles and cyclically rewriting
each individual cycle.  The order of a permutation is the least common
multiple of its disjoint cycle lengths.
\end{proposition}
\begin{proof}
Begin with a letter \(a\) and follow its successive images under \(\sigma\):
\[a,\ \sigma(a),\ \sigma^2(a),\ \ldots.\]
Because the set is finite and \(\sigma\)
is injective, the first repetition must be \(a\), giving one cycle.  Repeat
with a letter not yet used.  The resulting cycles have disjoint supports and
their product agrees with \(\sigma\) on every letter.  The orbit of each letter
under repeated application of \(\sigma\) determines its cycle, proving
uniqueness.  A power fixes every cycle exactly when its exponent is divisible
by that cycle length, so the least positive such exponent is the stated lcm.
\end{proof}

\begin{definition}
A \emph{transposition} is a two-cycle \((a\ b)\).  An \emph{adjacent
transposition} is one of the form \((i\ i+1)\).
\end{definition}

\begin{proposition}[Transpositions Generate \(S_n\)]
Every cycle is a product of transpositions, every permutation is a product of
transpositions, and every transposition is a product of adjacent transpositions.
\end{proposition}
\begin{proof}
Directly checking the images of the letters gives
\[(a_1\ a_2\ \cdots\ a_k)=(a_1\ a_k)(a_1\ a_{k-1})\cdots(a_1\ a_2).\]
The cycle-decomposition theorem therefore gives a transposition expression for
every permutation.  If \(i<j\), then
\[(i\ j)=(i\ i+1)\cdots(j-2\ j-1)(j-1\ j)(j-2\ j-1)\cdots(i\ i+1),\]
as is verified by following the images of \(i\), \(j\), and every other
letter.  Thus adjacent transpositions also generate \(S_n\).
\end{proof}

The preliminary definition of homomorphism introduced above is all that is
needed for the next result; its kernel and quotient theory are still deferred
to the later systematic treatment.

\begin{theorem}[Parity and the Sign Homomorphism]
Every expression of a fixed permutation as a product of transpositions has the
same parity of length.  Consequently
\[\operatorname{sgn}:S_n\longrightarrow\{1,-1\},\qquad
\operatorname{sgn}(\sigma)=(-1)^r\]
when \(\sigma\) is written as a product of \(r\) transpositions, is a
well-defined surjective homomorphism.
\end{theorem}
\begin{proof}
For \(\sigma\in S_n\), define its inversion number to be the number of pairs
\((i,j)\) with \(i<j\) and \(\sigma(i)>\sigma(j)\).  Multiplying by an
adjacent transposition on the right swaps two adjacent entries in the list
\(\sigma(1),\ldots,\sigma(n)\).  All inversions involving another entry are
unchanged in pairs, while the inversion between these two entries is reversed;
therefore the inversion number changes by exactly one.  The displayed
adjacent-transposition expression for \((i\ j)\) has
\(2(j-i)-1\) factors.  Thus every transposition changes inversion parity,
and a product of \(r\) transpositions has inversion number congruent to \(r\)
modulo two.  Two transposition expressions for the same permutation therefore
have the same parity.  The definition of \(\operatorname{sgn}\) is
well-defined.  Concatenating a transposition expression for \(\sigma\) with
one for \(\tau\) gives an expression for \(\sigma\tau\), so
\(\operatorname{sgn}(\sigma\tau)=\operatorname{sgn}(\sigma)
\operatorname{sgn}(\tau)\).  A transposition has sign \(-1\), so the map is
onto.
\end{proof}

\begin{definition}
The \emph{alternating group} \(A_n\) is the set of even permutations, namely
\(\{\sigma\in S_n:\operatorname{sgn}(\sigma)=1\}\).
\end{definition}

The sign rule shows directly that the product and inverse of even permutations
are even, so \(A_n\) is a subgroup.  The displayed sign rule says that
\(\operatorname{sgn}\) is a homomorphism.  The general theory of such maps
follows after cosets.  There we will prove that \(A_n\), as the kernel of the
sign map, is normal of index two.
For \(S_3\), \(A_3=\langle(123)\rangle\), and the three transpositions are
the odd permutations.  Later quotient theory will identify the two parity
classes with a copy of \(C_2\).

\section{Cosets, Index, and Lagrange's Theorem}

Every subgroup divides a group into pieces of equal size, whether or not its
left and right translates agree.  A coset is a translate of the subgroup: it remembers
how far an element is from belonging to \(H\), while forgetting the particular
element of \(H\) used to describe that displacement.

For \(H\leq G\), a \emph{left coset} is \(gH=\{gh:h\in H\}\); its number is
\([G:H]\), the \emph{index} of \(H\), when finite.

There are also right cosets \(Hg\).  They need not equal left cosets in a
nonabelian group.  A later condition, called normality, will make the two
notions coincide and permit a consistent multiplication of cosets.

\begin{lemma}
The left cosets of \(H\) partition \(G\), and each has the cardinality of
\(H\).
\end{lemma}
\begin{proof}
If \(g_1h_1=g_2h_2\), then \(g_2^{-1}g_1=h_2h_1^{-1}\in H\), which implies
\(g_1H=g_2H\).  Thus intersecting cosets are equal.  Every element lies in a
coset, and \(h\mapsto gh\) is a bijection \(H\to gH\).
\end{proof}

\begin{proposition}[Coset Equality Criterion]
For \(g,h\in G\), \(gH=hH\) if and only if \(h^{-1}g\in H\).
\end{proposition}
\begin{proof}
If \(gH=hH\), then \(g=ha\) for some \(a\in H\), so \(h^{-1}g\in H\).
Conversely, \(h^{-1}g=a\in H\) gives \(gH=haH=hH\).
\end{proof}

\begin{example}
The cosets of \(n\mathbb Z\) in \(\mathbb Z\) are the residue classes
\(r+n\mathbb Z\).  In \(S_3\), the left cosets of \(\langle(12)\rangle\) each
have two elements, so there are three of them.  This latter partition exists,
but its left and right cosets differ; later we will identify the precise
condition that permits their consistent multiplication.
\end{example}

\begin{example}
In \(D_4\), \(r\langle s\rangle=\{r,rs\}\), whereas
\(\langle s\rangle r=\{r,sr\}=\{r,r^{-1}s\}\).  The two sets differ, just as
a reflection and a rotation of a square do not commute.
\end{example}

\begin{theorem}[Lagrange]
If \(G\) is finite and \(H\leq G\), then \(|G|=[G:H]|H|\).  Consequently,
\(|g|\) divides \(|G|\) for every \(g\in G\).
\end{theorem}
\begin{proofidea}
The preceding lemma has already done the conceptual work.  Once \(G\) is a
disjoint union of translates of \(H\), counting the pieces gives the formula.
\end{proofidea}
\begin{proof}
Choose one representative from each left coset.  The preceding lemma says that
these cosets are pairwise disjoint, cover \(G\), and each has exactly \(|H|\)
elements.  If there are \([G:H]\) of them, counting the elements in this
disjoint union gives \(|G|=[G:H]|H|\).  Taking \(H=\langle g\rangle\) says
that the order of \(g\) divides \(|G|\).
\end{proof}

\begin{corollary}
If \(G\) is finite, then \(g^{|G|}=1\) for every \(g\in G\).
\end{corollary}
\begin{proof}
The order of \(g\) divides \(|G|\).
\end{proof}

\begin{corollary}
A finite group of prime order is cyclic.
\end{corollary}
\begin{proof}
Any element other than the identity has order dividing the prime group order,
so its order is the whole prime and it generates the group.
\end{proof}

The converse to Lagrange's theorem is false: divisibility of \(|H|\) by
\(|G|\) does not guarantee a subgroup of order \(|H|\).  For example, \(A_4\)
has no subgroup of order six, although six divides twelve.  This limitation is
one reason later existence theorems such as Cauchy's and Sylow's are valuable.

\section{Homomorphisms, Kernels, and Images}\label{sec:homomorphisms}

A homomorphism is the appropriate map between groups because it preserves the
operation.  Its kernel measures exactly what information the map discards,
and its image records exactly what it reaches.  These two subgroups are the
bridge between maps and quotient constructions.

\begin{definition}
\leavevmode\\
A \emph{homomorphism} \(\varphi:G\to K\) satisfies
\(\varphi(gh)=\varphi(g)\varphi(h)\).

Its \emph{kernel} is
\(\ker\varphi=\{g:\varphi(g)=1\}\), and its \emph{image} is
\(\operatorname{im}\varphi=\{\varphi(g):g\in G\}\).
\end{definition}

A homomorphism preserves the basic group operations automatically:
\(\varphi(1)=1\), \(\varphi(g^{-1})=\varphi(g)^{-1}\), and
\(\varphi(g^n)=\varphi(g)^n\).  Its kernel records exactly which elements the
map identifies, since \(\varphi(x)=\varphi(y)\) if and only if
\(x^{-1}y\in\ker\varphi\).  Images, by contrast, are subgroups but need not
be normal: \(\langle(12)\rangle\) is the image of \(C_2\to S_3\) but is not
invariant under conjugation.

For example, reduction modulo \(n\), \(\mathbb Z\to\mathbb Z/n\mathbb Z\),
has kernel \(n\mathbb Z\) and is onto.  The sign map
\(S_n\to\{\pm1\}\) has kernel \(A_n\).  The latter example is a useful source
of normal subgroups that is not obtained by inspection.  Likewise, the map
\(D_n\to\{\pm1\}\) that sends rotations to \(1\) and reflections to \(-1\)
has kernel \(\langle r\rangle\).

\begin{proposition}
The kernel and image of a homomorphism are subgroups; moreover,
\(\varphi\) is injective if and only if \(\ker\varphi=\{1\}\).
\end{proposition}
\begin{proof}
The subgroup test gives both subgroup assertions.  Also,
\(\varphi(x)=\varphi(y)\) exactly when
\(\varphi(xy^{-1})=1\), equivalently when \(xy^{-1}\in\ker\varphi\).
\end{proof}

\begin{definition}
A subgroup \(N\leq G\) is \emph{normal}, written \(N\trianglelefteq G\), if
\(gNg^{-1}=N\) for every \(g\in G\).
\end{definition}

\begin{proposition}[Equivalent Forms of Normality]
For \(N\leq G\), the conditions \(gNg^{-1}=N\) and \(gN=Ng\) for every
\(g\in G\) are equivalent.  Thus the left and right cosets of \(N\) coincide
exactly when \(N\) is normal.
\end{proposition}
\begin{proof}
Multiplying \(gNg^{-1}=N\) on the right by \(g\) gives \(gN=Ng\), and the
reverse implication follows by multiplying on the right by \(g^{-1}\).  The
coset statement is merely this equality written for each element \(g\).
\end{proof}

Normality says that left and right translates of \(N\) agree:
\(gN=Ng\) for every \(g\in G\).  Equivalently, membership in \(N\) is
unaffected by conjugation.  This condition is natural rather than cosmetic:
it is exactly what allows the product of two cosets to be independent of the
chosen representatives.

\begin{proposition}
\(\ker\varphi\trianglelefteq G\) for every homomorphism \(\varphi:G\to K\).
\end{proposition}
\begin{proof}
For \(n\in\ker\varphi\), \(\varphi(gng^{-1})=\varphi(g)1\varphi(g)^{-1}=1\).
The reverse inclusion follows on replacing \(g\) by \(g^{-1}\).
\end{proof}

\begin{corollary}
The alternating group \(A_n\) is a normal subgroup of \(S_n\) of index two.
\end{corollary}
\begin{proof}
The sign homomorphism has kernel \(A_n\), so the preceding proposition gives
normality.  It is surjective because a transposition has sign \(-1\).  Its
two fibers are the even and odd permutations; multiplication by a fixed
transposition pairs these two sets bijectively.  Thus each has \(n!/2\)
elements, and \([S_n:A_n]=2\).
\end{proof}

\section{Quotient Groups and Isomorphism Theorems}

The progression \(N\trianglelefteq G\Rightarrow G/N\) is fundamental.  We
identify two elements when they differ by an element of \(N\), then ask whether
the group operation descends to the resulting equivalence classes.  The
quotient criterion answers this question precisely.

\begin{theorem}[Quotient Criterion]
For \(N\leq G\), the rule \((gN)(hN)=(gh)N\) defines a group operation on
the left cosets if and only if \(N\trianglelefteq G\).  The group is \(G/N\).
\end{theorem}
\begin{proofidea}
Changing a representative of \(gN\) multiplies it by an element of \(N\).
Normality is exactly what lets that extra factor move past the representative
of the other coset and be absorbed back into \(N\).
\end{proofidea}
\begin{proof}
If \(N\) is normal, \(g'=gn_1\) and \(h'=hn_2\) imply
\[
g'h'=gh(h^{-1}n_1h)n_2\in ghN,
\]
so the rule is well-defined; associativity, identity \(N\), and inverse
\(g^{-1}N\) follow from \(G\).  Conversely, well-definedness of
\((gN)(nN)=(gN)N\) shows \(gng^{-1}\in N\) for every \(n\in N\), which is
normality.
\end{proof}

\begin{example}
The quotient \(\mathbb Z/n\mathbb Z\) identifies integers whose difference is
a multiple of \(n\).  The sign homomorphism has kernel \(A_n\), so
\(S_n/A_n\cong C_2\).  For the dihedral group
\(D_n=\langle r,s\mid r^n=s^2=1,\ srs^{-1}=r^{-1}\rangle\), the rotation
subgroup \(\langle r\rangle\) is normal and the quotient has two elements:
it records whether a symmetry preserves or reverses orientation.
\end{example}

\begin{example}[Why Normality Is Necessary]
In \(S_3\), let \(H=\langle(12)\rangle\).  It is not normal because
\((123)(12)(123)^{-1}=(23)\notin H\).  Indeed, the proposed coset product is
ambiguous: \(H=(12)H\), but
\[
H(13)H=(13)H\qquad\text{whereas}\qquad(12)H(13)H=(132)H,
\]
and these are different cosets.  The coset set exists, but it is not a
quotient group.
\end{example}

\begin{theorem}[Universal Property of a Quotient]
Let \(N\trianglelefteq G\), let \(\pi:G\to G/N\) be the quotient map, and
let \(f:G\to H\) be a homomorphism with \(N\subseteq\ker f\).  There is a
unique homomorphism \(\overline f:G/N\to H\) satisfying
\(f=\overline f\circ\pi\).
\end{theorem}
The following diagram gives a compact visual record of this factorization:
\[
\begin{tikzcd}[column sep=large, row sep=large]
G \arrow[r, "\pi"] \arrow[dr, "f"'] & G/N \arrow[d, "\overline f"] \\
& H
\end{tikzcd}
\]
A diagram \emph{commutes} when any two directed paths with the same source and
target define the same map.  Here commutativity says exactly
\(f=\overline f\circ\pi\).  This elementary language will be used whenever a
canonical map and a factorization clarify a later construction.
\begin{proofidea}
The quotient identifies \(g\) and \(gn\).  The condition \(f(n)=1\) says
that \(f\) also identifies them, so its value should depend only on the coset.
\end{proofidea}
\begin{proof}
Define \(\overline f(gN)=f(g)\).  If \(gN=g'N\), then \(g'=gn\) for some
\(n\in N\), and \(f(g')=f(g)f(n)=f(g)\); thus the definition is
well-defined.  It is a homomorphism because
\(\overline f((gN)(hN))=f(gh)=f(g)f(h)\).  The equation
\(f=\overline f\circ\pi\) holds by construction.  Finally every coset is
\(\pi(g)\), so this equation determines \(\overline f\) on every element of
\(G/N\), proving uniqueness.
\end{proof}

\begin{example}[A Factorization Computed in Full]
Let \(f:\mathbb Z\to C_6\) send \(1\) to a chosen generator \(u\).  Since
\(f(6k)=u^{6k}=1\), the subgroup \(6\mathbb Z\) is contained in \(\ker f\).
The universal property produces a unique map
\[\overline f:\mathbb Z/6\mathbb Z\longrightarrow C_6,
\qquad \overline f(\overline k)=u^k.\]
Why is this formula safe?  If \(\overline k=\overline\ell\), then
\(k-\ell=6q\), and \(u^k=u^{\ell+6q}=u^\ell(u^6)^q=u^\ell\).  The quotient
is not merely a convenient notation for residues: it is the unique group
through which every map that kills multiples of six must pass.
\end{example}

\begin{remark}[A Reusable Descent Test]
Whenever a proposed formula is defined on cosets, the first question is not
whether it is a homomorphism but whether it is well-defined.  The standard
test is: replace a representative by an equivalent one and show that the
output does not change.  In quotient groups this is usually equivalent to
showing that the subgroup being collapsed lies in a kernel.  The same sequence
of checks will be used for quotient rings and quotient modules.
\end{remark}

\begin{example}[Revisiting Cyclic Groups]
For \(\varphi_g:\mathbb Z\to G\), the image is \(\langle g\rangle\).  The
kernel calculation in the cyclic-group section therefore becomes the uniform
formula \(\langle g\rangle\cong\mathbb Z/\ker\varphi_g\).  When \(g\) has
infinite order this is \(\mathbb Z/\{0\}\cong\mathbb Z\), and when
\(|g|=n\) it is \(\mathbb Z/n\mathbb Z\).  This is precisely the later
theorem promised in the earlier direct proof.
\end{example}

\begin{theorem}[First Isomorphism Theorem]
For a homomorphism \(\varphi:G\to K\),
\[
G/\ker\varphi\cong\operatorname{im}\varphi.
\]
\end{theorem}
\begin{proofidea}
Two elements have the same image under \(\varphi\) precisely when their
quotient lies in \(\ker\varphi\).  Thus the fibers of \(\varphi\) are the
cosets of its kernel, so collapsing those cosets should leave exactly the
image.
\end{proofidea}
\begin{proof}
The map \(g\ker\varphi\mapsto\varphi(g)\) is well-defined because equal cosets
differ by an element of the kernel.  It is a surjective homomorphism with
trivial kernel, hence an isomorphism.
\end{proof}

\begin{remark}[How to Use the First Isomorphism Theorem]
The theorem is not a formula to apply after the fact.  It gives a workflow:
construct a homomorphism that records the feature of interest, calculate what
the homomorphism forgets, and then identify its image.  For the sign map the
feature is parity, the forgotten permutations are the even ones, and the image
is the two-element group.  For reduction modulo \(n\), the feature is a
residue class, the forgotten integers are the multiples of \(n\), and the
image is \(\mathbb Z/n\mathbb Z\).
\end{remark}

\begin{example}
The sign map \(\operatorname{sgn}:S_n\to\{\pm1\}\) is surjective and has
kernel \(A_n\).  The First Isomorphism Theorem therefore gives
\(S_n/A_n\cong\{\pm1\}\cong C_2\).  This is a typical use: calculate a
kernel, then recognize a quotient without multiplying cosets by hand.
\end{example}

\begin{theorem}[Second Isomorphism Theorem]
Let \(N\trianglelefteq G\) and \(H\leq G\).  Then \(HN\leq G\),
\(H\cap N\trianglelefteq H\), and
\[HN/N\cong H/(H\cap N).\]
\end{theorem}
\begin{proof}
The product \(HN\) is a subgroup because normality of \(N\) gives
\[
  (h_1n_1)(h_2n_2)=h_1h_2(h_2^{-1}n_1h_2)n_2\in HN.
\]
The subgroup test gives the remaining closure.  Conjugating an element
of \(H\cap N\) by an element of \(H\) keeps it in both \(H\) and \(N\), so it
is normal in \(H\).  Restrict the quotient map \(G\to G/N\) to \(H\).  Its
kernel is \(H\cap N\), and its image is \(HN/N\).  The First Isomorphism
Theorem gives the claimed isomorphism.
\end{proof}

The theorem says that when \(H\) enters the quotient by \(N\), precisely the
elements already in \(N\), namely \(H\cap N\), are killed.

\begin{theorem}[Third Isomorphism Theorem]
If \(N\trianglelefteq G\) and \(N\leq H\trianglelefteq G\), then
\(H/N\trianglelefteq G/N\) and \((G/N)/(H/N)\cong G/H\).
\end{theorem}
\begin{proof}
Define \(\psi:G/N\to G/H\) by \(\psi(gN)=gH\).  If \(gN=g'N\), then
\(g^{-1}g'\in N\subseteq H\), so \(gH=g'H\); hence \(\psi\) is well-defined.
It is a surjective homomorphism, and its kernel consists of the cosets \(gN\)
with \(g\in H\), which is \(H/N\).  The First Isomorphism Theorem proves the
claim, including normality of the kernel.
\end{proof}

\begin{example}[Quotienting in Stages]
The chain \(12\mathbb Z\subseteq4\mathbb Z\subseteq\mathbb Z\) yields
\[(\mathbb Z/12\mathbb Z)/(4\mathbb Z/12\mathbb Z)
\cong\mathbb Z/4\mathbb Z.\]
The left side first remembers residues modulo twelve and then identifies
residues differing by four; the right side simply remembers residues modulo
four.  This is precisely the content of quotienting in stages.
\end{example}

This is the slogan \highlight{quotienting in stages is the same as quotienting
once}.  It is a basic organizational principle in all later quotient theories.

\begin{theorem}[Correspondence Theorem]
Let \(N\trianglelefteq G\), and let \(\pi:G\to G/N\) be the quotient map.
Then \(H\mapsto H/N\) is an inclusion-preserving bijection between subgroups
of \(G\) containing \(N\) and subgroups of \(G/N\).  Its inverse is
\(K\mapsto\pi^{-1}(K)\).  Under this bijection, normal subgroups correspond,
and \(H\trianglelefteq G\) implies \((G/N)/(H/N)\cong G/H\).
\end{theorem}
\begin{proof}
If \(N\leq H\), then \(H/N\leq G/N\).  If \(K\leq G/N\), then
\(\pi^{-1}(K)\leq G\) and contains \(N\).  The two operations are inverse
because \(\pi\) is surjective: \(\pi^{-1}(H/N)=H\) and
\(\pi(\pi^{-1}(K))=K\).  Conjugation commutes with \(\pi\), proving the
normality assertion.  The final statement is the Third Isomorphism Theorem.
\end{proof}

For instance, subgroups of \(\mathbb Z/n\mathbb Z\) correspond to subgroups
of \(\mathbb Z\) that contain \(n\mathbb Z\).  The containment condition is
necessary: only such subgroups are unions of cosets of \(n\mathbb Z\), so only
they have a meaningful image as a subgroup of the quotient.

\subsection*{Worked Quotient Examples}

The following two examples gather the ideas of normality, quotient groups,
universal properties, and isomorphism theorems in concrete settings.

\begin{example}[Reduction Modulo \(12\)]
Consider \(\pi:\mathbb Z\to\mathbb Z/12\mathbb Z\).  Two integers have the
same image exactly when their difference is a multiple of twelve:
\[
\pi(a)=\pi(b)\quad\Longleftrightarrow\quad a-b\in12\mathbb Z.
\]
Thus the quotient identifies exactly the elements that the map cannot
distinguish.  The subgroup correspondence says that subgroups of
\(\mathbb Z/12\mathbb Z\) correspond to the subgroups of \(\mathbb Z\)
containing \(12\mathbb Z\).  For instance, the subgroup generated by
\(\overline3\) has four elements and corresponds to \(3\mathbb Z\), while the
subgroup generated by \(\overline4\) has three elements and corresponds to
\(4\mathbb Z\).  Inclusion reverses neither upstairs nor downstairs: the
containment \(12\mathbb Z\subseteq4\mathbb Z\subseteq\mathbb Z\) becomes
\(\{\overline0\}\subseteq\langle\overline4\rangle\subseteq
\mathbb Z/12\mathbb Z\).
\end{example}

\begin{example}[The Sign Quotient of \(S_3\)]
The sign map separates the six permutations of \(S_3\) into the three even
permutations \(A_3=\{1,(123),(132)\}\) and the three odd permutations.  The
two cosets of \(A_3\) are therefore \(A_3\) itself and \((12)A_3\).  In the
quotient, every even permutation becomes the identity coset and every odd
permutation becomes the same nontrivial coset.  Its square is the identity,
because the product of two odd permutations is even.  Hence the quotient is
the two-element group.  This makes the abstract statement
\(S_3/A_3\cong C_2\) concrete: the quotient retains parity and forgets every
other feature of a permutation.
\end{example}

\begin{proposition}[A Direct Product of Cyclic Groups]
For positive integers \(m,n\), the group \(C_m\times C_n\) is cyclic if and
only if \(\gcd(m,n)=1\).  In that case \((g,h)\), where \(g\) and \(h\) are
generators of the two factors, generates the product.
\end{proposition}
\begin{proof}
The order of \((g,h)\) is the least positive integer \(k\) for which both
\(g^k=1\) and \(h^k=1\).  Thus its order is \(\operatorname{lcm}(m,n)\).
If \(m\) and \(n\) are coprime, this least common multiple is \(mn\), the
order of the product, so \((g,h)\) is a generator.  Conversely, suppose a
single element generated \(C_m\times C_n\).  Its order would be \(mn\), but
the order of every element divides \(\operatorname{lcm}(m,n)\).  Therefore
\(mn\mid\operatorname{lcm}(m,n)\), which forces \(\gcd(m,n)=1\).
\end{proof}

This example contrasts direct products with quotients.  A quotient imposes
identifications and thereby discards distinctions; a direct product places two
independent systems side by side.  Their universal properties will recur
throughout algebra, but their concrete group-theoretic forms already show two
basic operations of algebraic construction.

\subsection*{A Guided Tour of the Quotient Construction}

It is useful to see the quotient construction operate several times before
leaving Chapter~2.  The same four questions organize each instance:

\begin{enumerate}
\item What elements should be regarded as indistinguishable?
\item Do these equivalence classes form cosets of a normal subgroup?
\item Does the proposed operation depend only on the classes, rather than on
      chosen representatives?
\item What information is retained by the resulting quotient?
\end{enumerate}

For \(\mathbb Z/n\mathbb Z\), the answer to the first question is that two
integers are indistinguishable when their difference is a multiple of \(n\).
The retained information is a remainder.  For \(S_3/A_3\), two permutations are
indistinguishable when they have the same parity; the retained information is
whether the permutation is even or odd.  For \(D_n/\langle r\rangle\), two
symmetries are indistinguishable when they differ by a rotation; the quotient
retains whether the symmetry preserves or reverses orientation.

\begin{example}[The Dihedral Quotient in Detail]
Let \(N=\langle r\rangle\leq D_n\).  The relation \(srs^{-1}=r^{-1}\)
implies \(sr^is^{-1}=r^{-i}\in N\), so conjugation by the reflection keeps
the rotation subgroup inside itself.  Conjugation by a rotation plainly does
the same; therefore \(N\) is normal.  Its two cosets are
\[N=\{1,r,\ldots,r^{n-1}\}\quad\text{and}\quad sN=
\{s,sr,\ldots,sr^{n-1}\}.\]
The product of two elements of the second coset lies in \(N\), because
\[(sr^i)(sr^j)=s(r^is)r^j=r^{j-i}.\]
Thus the nontrivial coset has order two
in the quotient.  This gives a direct, representative-level verification that
\(D_n/N\cong C_2\).
\end{example}

\begin{example}[A Nonnormal Subgroup Revisited]
Let \(H=\langle(12)\rangle\leq S_3\).  It has three left cosets, each with two
elements.  The set of cosets is therefore perfectly meaningful, but the
attempted product of cosets depends on representatives.  Indeed, \(H=(12)H\),
yet multiplying each representative by the coset \((13)H\) produces different
cosets.  The obstruction is not a flaw in the formula; it is the fact that
conjugating \((12)\) by \((13)\) produces a transposition outside \(H\).
The normality condition is exactly what removes this obstruction in the
quotient criterion.
\end{example}

\begin{remark}[Proof Technique: Pass to the Right Quotient]
When a problem contains a homomorphism, a normal subgroup, or a statement that
is unchanged after ``forgetting'' some information, try forming a quotient.
The quotient may turn a complicated relation into the identity relation.  The
First Isomorphism Theorem then identifies the simplified object with an image,
and the Correspondence Theorem translates subgroup questions upstairs and
downstairs.  This is one of the central reusable moves of graduate algebra.
\end{remark}

\subsection*{Why These Constructions Matter Later}

The material of this chapter forms a prototype for several later branches of
algebra.  A quotient group is the first example of imposing an equivalence
relation compatible with an operation.  In ring theory, ideals play the role
of normal subgroups, and the same well-definedness question produces quotient
rings.  In module theory, submodules replace normal subgroups; the isomorphism
theorems and correspondence theorem reappear with almost identical proofs.

Group actions supply a second durable pattern.  When a group acts on a set of
geometric objects, on roots of a polynomial, or on a vector space, the orbit
records what can be moved into what, and the stabilizer records the residual
symmetry at a chosen point.  This is why the elementary formula
\(|G|=|Gx|\,|G_x|\) later becomes a tool in Galois theory, representation
theory, algebraic geometry, and number theory.

Finally, the examples \(S_3\) and \(D_4\) should be regarded as small models
of two different sources of noncommutativity.  In \(S_3\), noncommutativity
comes from rearranging different objects; in \(D_4\), it comes from composing
rotations with reflections.  Both groups have meaningful quotients and
semidirect-product decompositions.  Revisiting them when later constructions
become abstract is an efficient way to keep the structural ideas grounded.

\section{Direct Products and Extensions}

The external direct product \(G\times H\) combines independent symmetries:
one may act in the first coordinate without affecting the second.  For example,
the symmetries of a rectangle with unequal side lengths form \(C_2\times C_2\).
The internal criterion below recognizes this construction inside a larger group.

\begin{proposition}
Suppose \(N,H\trianglelefteq G\), \(G=NH\), and \(N\cap H=\{1\}\).  Then
\(G\cong N\times H\).
\end{proposition}
\begin{proof}
Every element is \(nh\).  Such an expression is unique because
\(n_1h_1=n_2h_2\) implies \(n_2^{-1}n_1=h_2h_1^{-1}\in N\cap H\).
Normality makes the commutator of an element of \(N\) with one of \(H\) lie in
both subgroups, hence be trivial.  Therefore \((n,h)\mapsto nh\) is a
bijective homomorphism.
\end{proof}

\begin{remark}[Looking Ahead]
The quotient map \(G\to G/N\) is the prototype for quotient rings and quotient
modules.  Kernels, images, and quotients will later be packaged into exact
sequences.
\end{remark}

\section*{Exercises}
\begin{exercise}[Basic]
In the multiplicative group \(\mathbb Q^\times\), prove that
\[\langle-1,2\rangle=\{\pm2^k:k\in\mathbb Z\}.\]
Decide whether this subgroup is cyclic.
\end{exercise}
\begin{exercise}[Core]
Prove that a subgroup of index two is normal.
\end{exercise}
\begin{exercise}[Core]
Let \(C_{12}=\langle g\rangle\).  Determine
\(\langle g^3\rangle\cap\langle g^4\rangle\) and the subgroup generated by
\(\langle g^3\rangle\cup\langle g^4\rangle\).  State the corresponding
gcd/lcm pattern for two subgroups of a finite cyclic group.
\end{exercise}
\begin{exercise}[Core]
Let \(H\leq G\) be the unique subgroup of \(G\) having its order.  Prove
that \(H\trianglelefteq G\).  Apply this criterion to the subgroup of order
three in a cyclic group of order twelve.
\end{exercise}
\begin{exercise}[Core]
Let \(N\trianglelefteq G\), and let \(gN\in G/N\).  Prove that the order of
\(gN\) is the least positive integer \(r\) for which \(g^r\in N\), when
such an integer exists.  Apply this to compute the order of \(r\langle r^d\rangle\)
in \(D_n/\langle r^d\rangle\) for \(d\mid n\).
\end{exercise}
\begin{exercise}[Further]
If \(G\) is cyclic of order \(n\), prove
\(\operatorname{Aut}(G)\cong(\mathbb Z/n\mathbb Z)^\times\).
\end{exercise}
\begin{exercise}[Looking Ahead]
Let \(N\trianglelefteq G\).  Use the Correspondence Theorem to show that
maximal subgroups of \(G/N\) correspond to maximal subgroups of \(G\) that
contain \(N\).
\end{exercise}
